\chapter*{Ход выполнения практики}
\addcontentsline{toc}{chapter}{Ход выполнения практики}

\section*{Задания и отметки о выполнении}

\task{Изучить технологию OpenMP и показатели эффективности параллельных
программ для вычислительных систем с общей памятью.}

\task{Разработать параллельную OpenMP-версию расчётного модуля программного
комплекса HeatSim3D с управляемым числом потоков.}

\task{Обеспечить сборку и работу программного комплекса в операционных
системах Linux и Windows.}

\task{Проверить корректность параллельной версии путём сравнения
с аналитическим решением и с результатами последовательного расчёта.}

\task{Исследовать ускорение и эффективность параллельной версии
для расчётных сеток разного размера.}

\task{Подготовить отчёт по практике в соответствии с требованиями
нормоконтроля кафедры информационных систем и компьютерного
моделирования.}

\clearpage
\section*{Отзыв о работе обучающегося}

\fillme{Отзыв пишет руководитель практики. Ниже --- черновик для согласования.}

В ходе практики Богунь М. В. разработал параллельную версию программного
комплекса HeatSim3D для численного моделирования нестационарной
теплопроводности с использованием технологии OpenMP. Студент провёл анализ
зависимостей по данным в алгоритме явной разностной схемы, обеспечил работу
программы в операционных системах Linux и Windows, выполнил проверку
корректности по аналитическому решению и исследовал ускорение и
эффективность параллельной версии на сетках разного размера.

Студент продемонстрировал умение планировать вычислительный эксперимент,
количественно оценивать производительность программ и объяснять
полученные результаты. Задачи практики выполнены в полном объёме.

\progressapprov
